\subsection{A relative-trace obstruction without signed enumeration}
\begin{proposition}[Relative-trace criterion]
\label{prop:relative-trace}
Let $R$ be a commutative fusion ring with self-dual basis, and put
$B=\sum_{i,j,k}(N_{ij}^k)^2$. Suppose $L\in\Z_{>0}$ and $f$ are
formal codegrees of $R$, and set $F=\mathbb Q(f)$.
If every nonzero totally positive element of
$L\mathcal O_F\cap f\mathcal O_F$ has field trace greater than $B$,
then $R$ has no categorification over $\C$.
\end{proposition}
\begin{proof}
Sphericalize a hypothetical category and take its quotient by
$\varepsilon=-1$. Both lifts of a self-dual simple remain self-dual
\cite[Section~5.1(a)]{OstrikRank3}; hence the signed multiplication
matrices are symmetric and commute. Writing them as $M_i$, their
Casimir satisfies
\[
 \operatorname{Tr}\Bigl(\sum_iM_i^2\Bigr)
 =\sum_{i,j,k}(T_{ij}^k)^2\le B.
\]
Let $d_i$ be the lift dimensions, $K=\mathbb Q(d_i)$, and
$\delta=\sum_i d_i^2$. Every embedding of $K$ gives a distinct
dimension character and contributes its conjugate of $\delta$ as
a Casimir eigenvalue. This counts $[K:\mathbb Q]$ characters even
when some eigenvalues coincide. All Casimir eigenvalues are positive,
so $\operatorname{Tr}_{K/\mathbb Q}(\delta)\le B$.

The sphericalization has dimension $2\delta$ and formal codegrees
$2L,2f$. The induced-unit theorem gives algebraic integers
$\delta/L,\delta/f$ in its dimension field $K$
\cite[Theorem~2.13 and Corollaries~2.14--2.15]{OstrikRank3}.
Thus $F\subseteq K$. The element
$t=\operatorname{Tr}_{K/F}(\delta)$ is totally positive and lies in
$L\mathcal O_F\cap f\mathcal O_F$. Its trace is at most $B$, a
contradiction. No equality $K=\mathbb Q(\delta)$ is assumed.
\end{proof}

\begin{proposition}
The rings $R7\_4\_1417$ and $R7\_4\_1420$ admit no complex
fusion-category categorification.
\end{proposition}
\begin{proof}
Both rings have formal codegrees $20$ and
$f=510+70\sqrt{51}$, whose minimal polynomial is
$z^2-1020z+10200$. Here
$\mathcal O_F=\Z[\sqrt{51}]$. Write
$t=20(x+y\sqrt{51})$ in the intersection ideal. Then
\[
 \frac{t}{f}=(x-7y)+\left(y-\frac{7x}{51}\right)\sqrt{51}.
\]
Integrality forces $51\mid x$, while total positivity forces $x>0$.
Consequently $\operatorname{Tr}_{F/\mathbb Q}(t)=40x\ge2040$.
The exact Casimir trace bounds are $1086$ and $1070$, respectively,
so Proposition~\ref{prop:relative-trace} applies. The source-bound
checker recomputes both characteristic polynomials, trace bounds,
and the displayed quotient identity using exact arithmetic.
\end{proof}
